% Intended LaTeX compiler: pdflatex
\documentclass[a4paper,12pt]{article}
\usepackage[utf8]{inputenc}
\usepackage[T1]{fontenc}
\usepackage{graphicx}
\usepackage{longtable}
\usepackage{wrapfig}
\usepackage{rotating}
\usepackage[normalem]{ulem}
\usepackage{amsmath}
\usepackage{amssymb}
\usepackage{capt-of}
\usepackage{hyperref}
\usepackage[margin=0.2in]{geometry}
\usepackage{amsmath}
\usepackage{amsfonts}
\usepackage{indentfirst}
\usepackage{pgffor}
\usepackage{amssymb}
\usepackage{cancel}
\usepackage{amsthm}
\usepackage{polynom}
\usepackage{tikz}
\usepackage[tikz]{bclogo}
\usepackage{listings}
\usepackage[most]{tcolorbox}
\usepackage{forest}
\usepackage{adjustbox}
\usepackage{tikz-3dplot}
\usepackage{mathtools}
\usepackage{pgfplots}
\usetikzlibrary{tikzmark,calc,fit,matrix,arrows,automata,positioning}
\usepackage{centernot}
\usepackage{lmodern}
\usepackage{physics}
\usepackage[boxed,linesnumbered,vlined]{algorithm2e}
\usepackage{algpseudocode}
\usepackage{algorithmicx}
\usepackage{svg}
\tikzstyle{every picture}+=[remember picture]
\DeclareMathOperator{\spn}{span}
\DeclareMathOperator{\dom}{domain}
\DeclareMathOperator{\ran}{range}
\DeclareMathOperator{\rng}{range}
\DeclareMathOperator{\img}{Im}
\DeclareMathOperator{\adj}{adj}
\newcommand\dif[1]{\:\textrm{d}#1}
\DeclarePairedDelimiter\ceil{\lceil}{\rceil}
\DeclarePairedDelimiter\floor{\lfloor}{\rfloor}
\newcommand{\ihat}{\hat{\textbf{\i}}}
\newcommand{\jhat}{\hat{\textbf{\j}}}
\newcommand{\khat}{\hat{\textbf{k}}}
\newcommand{\rhat}{\hat{\textbf{r}}}
\newcommand{\thetahat}{\boldsymbol{\hat{\theta}}}
\author{mahmood}
\date{\today}
\title{binary coding}
\hypersetup{
 pdfauthor={mahmood},
 pdftitle={binary coding},
 pdfkeywords={},
 pdfsubject={},
 pdfcreator={Emacs 30.0.50 (Org mode 9.6)}, 
 pdflang={English}}
\begin{document}

\maketitle
\definecolor{Brown}{cmyk}{0,0.81,1,0.60}
\definecolor{OliveGreen}{cmyk}{0.64,0,0.95,0.40}
\definecolor{CadetBlue}{cmyk}{0.62,0.57,0.23,0}
\definecolor{lightlightgray}{gray}{0.9}
\lstset{
  language=C,                             % Code langugage
  basicstyle=\ttfamily,                   % Code font, Examples: \footnotesize, \ttfamily
  keywordstyle=\color{OliveGreen},        % Keywords font ('*' = uppercase)
  commentstyle=\color{gray},              % Comments font
  numbers=left,                           % Line nums position
  numberstyle=\tiny,                      % Line-numbers fonts
  stepnumber=1,                           % Step between two line-numbers
  numbersep=5pt,                          % How far are line-numbers from code
  backgroundcolor=\color{lightlightgray}, % Choose background color
  frame=single,                           % A frame around the code
  tabsize=2,                              % Default tab size
  captionpos=b,                           % Caption-position = bottom
  breaklines=true,                        % Automatic line breaking?
  breakatwhitespace=false,                % Automatic breaks only at whitespace?
  showspaces=false,                       % Dont make spaces visible
  showtabs=false,                         % Dont make tabls visible
  columns=flexible,                       % Column format
  morekeywords={__global__, __device__},  % CUDA specific keywords
  showstringspaces=false,                 % dont show spaces in strings
}
% so that latex doesnt complain about sage not being a language
\lstdefinelanguage{sage}{}

% environment definition for special blocks
\theoremstyle{definition}
\newtheorem{my_example}{Example}
\counterwithin{my_example}{section}
\counterwithin{my_example}{subsection}
\newtheorem{definition}{Definition}
\counterwithin{definition}{section}
\counterwithin{definition}{subsection}
\newtheorem{characteristic}{Characteristic}
\newtheorem{entailment}{Entailment}
%\newtheore*{proof}{Proof}
\newtheorem{lemma}{Lemma}
\newtheorem{question}{Question}
\newtheorem{subquestion}{Subquestion}[question]
%\newtheorem{note}{Note}
\newtheorem{answer}{Answer}
\counterwithin{answer}{question}
\counterwithin{answer}{subquestion}
\newtheorem{step}{Step}[answer]

% the following snippet auto resizes images to fit properly
% Determine if the image is too wide for the page.
% \makeatletter
% \def\ScaleIfNeeded{%
%   \ifdim\Gin@nat@width>\linewidth
%     \linewidth
%   \else
%     \Gin@nat@width
%   \fi
% }
% \makeatother
% % Resize figures that are too wide for the page.
% \let\oldincludegraphics\includegraphics
% \renewcommand\includegraphics[2][]{%
%   \oldincludegraphics[width=\ScaleIfNeeded]{#2}
% }

\newenvironment{note}
  {\begin{bclogo}[logo=\bcinfo, couleurBarre=orange, couleur=lightgray]{ note}}
  {\end{bclogo}}
\newenvironment{attention}
  {\begin{bclogo}[logo=\bcattention, couleurBarre=red, noborder=true, 
                couleur=red!15]{Important!}}
  {\end{bclogo}}

\begin{definition}
a \textbf{binary coding} is an \href{discrete_maths2.org}{injective} \href{discrete_maths2.org}{function} that maps each character of a finite set of characters to string of bits, given a character \(c\), and a binary coding \(k\), \(k(c)\) is called the \textbf{coding} of \(c\), and the number of bits in \(k(c)\) is denoted by \(|k(c)|\)
\begin{my_example}
consider the following binary coding function for the set \(C=\{a,b,c,d,e,f\}\)
\begin{center}
\begin{tabular}{lr}
character \(c\) & coding \(k(c)\)\\[0pt]
\hline
a & 000\\[0pt]
b & 001\\[0pt]
c & 010\\[0pt]
d & 011\\[0pt]
e & 100\\[0pt]
f & 101\\[0pt]
\end{tabular}
\end{center}
\end{my_example}
\end{definition}
\begin{definition}
the binary coding of a \href{20230119162026-file.org}{file} \(F=c_1,c_2,\dots,c_l\) is a chain of codes of its characters, \(k(F)=k(c_1)\circ\dots\circ k(c_1)\)
\begin{my_example}
\[
k_1(adf) = 000 \circ 011 \circ 101 = 000011101
\]
\end{my_example}
\end{definition}
\#+begin\textsubscript{definition}
the \textbf{frequency} of a character \(c \in C\) in a file \(F\), \(f(c,F)\) is the number of times that character \(c\) appears in \(F\)
\begin{characteristic}
let \(|k(F)|\) denote the number of characters in an arbitrary file \(F\)
$\backslash$[
\begin{center}
\begin{tabular}{llll}
k(F) & =\(\sum\)\textsubscript{c \(\in\) C} & k(c) & f(c,F)\\[0pt]
\end{tabular}
\end{center}
$\backslash$]
\end{characteristic}
\begin{my_example}
suppose to we wanna encode a file \(F\) that contains 100 words over a 6-letter alphabet, we also assume that each character in \(F\) has a different frequency
an efficient coding of \(F\) would match short binary strings with characters that have a high frequency, and long strings wirth letters that have lower frequency, the following table compares two codings of \(F\):
\begin{center}
\begin{tabular}{lrrr}
character & frequency & \(k_1\), \href{20230119170007-fixed_length_code_word.org}{fixed-length code word} & \(k_2\), \href{20230119170058-variable_length_code_word.org}{variable-length code word}\\[0pt]
\hline
a & 45 & 000 & 0\\[0pt]
b & 13 & 001 & 101\\[0pt]
c & 12 & 010 & 100\\[0pt]
d & 16 & 011 & 111\\[0pt]
e & 9 & 100 & 1101\\[0pt]
f & 5 & 101 & 1100\\[0pt]
\end{tabular}
\end{center}
with this example table we can already see why variable length code is more efficient as \(k_2(F)=222,k_1(F)=300\)
\end{my_example}
\end{document}